% Generado por Code/tools/tablas_latex.py desde 8_spearman.csv. No editar a mano: volver a correr el script.
\begin{table}[H]
\centering
\small
\caption{Correlación de Spearman entre el orden de separabilidad y el de clasificación de las cinco condiciones (M8).}
\label{tab:m8_spearman}
\begin{adjustbox}{max width=\linewidth}
\begin{tabular}{l l r r}
\toprule
Subconjunto & Métrica de separabilidad & $\rho_s$ & $p$ exacto bilateral \\
\midrule
Masas & Davies--Bouldin & 0.718 & 0.167 \\
Masas & Fisher $J$ & \textbf{0.821} & 0.133 \\
Masas & KTA lineal & 0.718 & 0.167 \\
Calcificaciones & Davies--Bouldin & \textbf{0.975} & 0.033 \\
Calcificaciones & Fisher $J$ & 0.667 & 0.267 \\
Calcificaciones & KTA lineal & \textbf{0.975} & 0.033 \\
\bottomrule
\end{tabular}
\end{adjustbox}
\par\smallskip
\parbox{0.95\linewidth}{\footnotesize Submuestra de 200 lesiones frente al AUC de prueba, $k = 12$. Con $n = 5$ hay 120 órdenes; $\rho_s$ es la correlación de Pearson de los rangos promedio, porque C1 y C2 empatan. Es descriptivo: con el empate, el mejor orden alcanzable da $p = 4/120$. En negritas, la métrica con mayor $\rho_s$ en cada subconjunto.}
\end{table}
